\section{According to His Works, and What May Not Be Said}\label{sec:works}

Two matters remain, and both are places where the doctrine of judgment is regularly pushed past its evidence---once in the direction of anxiety and once in the direction of presumption.

\subsection{The criterion, and what it is not a claim about}

One phrase runs through every witness examined in this article. Matthew: ``then will he render to every man according to his works.'' Romans: ``Who will render to every man according to his works.'' The Apocalypse, twice in three verses: ``the dead were judged by those things which were written in the books, according to their works,'' and ``they were judged, every one according to their works.'' Paul at Corinth: ``that every one may receive the proper things of the body, according as he hath done, whether it be good or evil.'' Lateran~IV: \emph{redditurus singulis secundum opera sua}. \emph{Benedictus Deus} quotes the Corinthian clause in the very sentence in which it retains the general judgment. And Matthew's judgment scene supplies six works and no other criterion at all.

The uniformity is not decorative. Whatever else the last day is, it is an assessment of what a man did.

It is also, and this needs saying in the same breath, not a statement about what causes a man to be saved. The distinction is not a modern nervousness; it is where Trent placed the emphasis in the same session that anathematized the denial of temporal punishment. Canon~32 anathematizes anyone who says that the good works of a justified man are the gift of God \emph{in such a way} that they are not also his own good merits---and then, in the same sentence, describes those works as the ones ``which he performs through the grace of God and the merit of Jesus Christ, whose living member he is.''\endnote{Council of Trent, session~VI, canon~32, Waterworth's English, read 25 July 2026 at the registered page images of printed pp.~48--49 of the digitized Burns and Oates reprint; the canon runs across the page break and both leaves were inspected. The Latin at the registered page image of printed p.~38 of the Tauchnitz reprint reads \emph{aut ipsum iustificatum bonis operibus, quae ab eo per Dei gratiam et Iesu Christi meritum, cuius vivum membrum est, fiunt, non vere mereri augmentum gratiae, vitam aeternam \ldots}. The canon's conclusion, ``if so be, however, that he depart in grace,'' is on the second leaf.} The merit is real and it is the man's; the whole of its power is grace and the merit of Christ, of whom he is a living member. And the very next canon exists to prevent the doctrine from being read as an encroachment: anyone who says that by this doctrine of justification ``the glory of God, or the merits of our Lord Jesus Christ are in any way derogated from'' is anathema.\endnote{Council of Trent, session~VI, canon~33, Waterworth's English, read at the registered page image of printed p.~49. The full text of the canon appears entire on that leaf.}

This article does not develop the doctrine of justification; that treatise is not its subject and cannot be compressed without distorting it. What belongs here is only the boundary. ``According to works'' is a statement about the \emph{measure} of the judgment, not about the \emph{source} of the standing being measured. The six works of Matthew~25 are not a schedule of purchases. They are the shape a life takes when charity is in it, which is why the just in the scene do not recognize their own record: charity does not keep accounts of itself, and a man who could produce the ledger would be exhibiting something other than charity.\endnote{The reading of the two surprised questions in Matthew 25:37--39 and 25:44 offered in this sentence is this article's own and is labelled as project synthesis in the terminal appendix. That both groups ask the same question in nearly the same words is a fact of the text; what it signifies is not stated in the text.}

\subsection{The time, which is not merely undefined but withheld}

On the date the tradition is not merely silent but instructed. ``But of that day and hour no one knoweth: no, not the angels of heaven, but the Father alone.'' The Catechism repeats it in its own voice: ``The Last Judgment will come when Christ returns in glory. Only the Father knows the day and the hour; only he determines the moment of its coming.''\endnote{\emph{Catechism of the Catholic Church} 1040, official English, read 25 July 2026 at the registered archived page. Matthew 24:36 from the tracked Challoner Douay--Rheims, with the American-edition variant recorded at its first occurrence in Section~\ref{sec:defined}.}

The medieval treatment asks the question formally, and its answer turns out to be the same argument that carried Section~\ref{sec:trial}. God communicates to creatures both causal power and knowledge, but reserves to himself both the operations in which no creature co-operates and the knowledge of them. The end of the world is exactly such an operation: ``the world will come to an end by no created cause, even as it derived its existence immediately from God. Wherefore the knowledge of the end of the world is fittingly reserved to God''---and the text offers as the Lord's own statement of the reason, ``It is not for you to know the times or moments which the Father hath put in His own power.''\endnote{\emph{Supplementum} q.~88, a.~3, corpus, English Dominican translation read 25 July 2026 at the registered New Advent page; the article rubric \emph{Utrum tempus futuri iudicii sit ignotum} and its parallel-place citations were read at the registered Leonine page image of printed p.~205. The scriptural tag is Acts 1:7. Note that this is the same premise as q.~88, a.~1: the general judgment corresponds to the operation by which all things proceeded immediately from God, and immediacy to God is what makes both the event and its date God's alone. The reply to objection~4 adds the pastoral reason, that the uncertainty conduces to watchfulness in two ways, the second being that a man must be careful not only of himself but of his family, city, kingdom, or the whole Church, whose duration is not measured by his own. Reply~2 is a sustained refusal, on Augustine's authority, to read the gospel signs as a calendar: ``it is impossible to decide after how long a time it will take place, nor fix the month, year, century, or thousand years.''} That is why every attempt to date the last judgment is not merely mistaken but a trespass on something reserved; and it fails in a further way worth noticing: since the particular judgment is at death and no one knows the hour of that either, a correct date for the end of the world would be of no use to anybody. The two ignorances are the same ignorance in two forms, and the response to both is the one the eschatological discourse itself prescribes: ``Watch ye therefore, because you know not the day nor the hour.''

\subsection{No judgment about any person}

The last restriction is the hardest to keep and the most important.

Paul states it as a rule of Christian conduct and gives the reason: ``Therefore, judge not before the time: until the Lord come, who both will bring to light the hidden things of darkness and will make manifest the counsels of the hearts.'' The prohibition is not on discernment but on \emph{verdicts}, and the ground is chronological rather than sentimental: the evidence is not in, and will not be in until the Lord comes. Paul adds the most uncomfortable clause in the passage, and it cuts against self-acquittal as much as against the condemnation of others: ``For I am not conscious to myself of anything. Yet am I not hereby justified: but he that judgeth me is the Lord.''

The Catechism states the rule in its own terms, in the middle of its treatment of mortal sin, where it is least convenient and most needed: ``although we can judge that an act is in itself a grave offense, we must entrust judgment of persons to the justice and mercy of God.''\endnote{\emph{Catechism of the Catholic Church} 1861, official English, read 25 July 2026 at the registered archived page. The sentence stands at the end of the paragraph on mortal sin as ``a radical possibility of human freedom.'' The distinction it draws---between judging an act and judging a person---is the operative one for everything in this subsection.} Acts have a moral species that can be known; persons have an interior history that cannot. And the doctrine examined in this article is the reason: the judgment that reveals the counsels of hearts has not happened yet, and it is reserved to one judge who is not any of us.

This bears directly on how the material in the preceding sections may be used. The scene in Matthew~25 is a description of a judgment Christ will make; it is not a diagnostic for use on one's neighbours or one's relatives or oneself. The definitions of Lyons, Benedict~XII, and Florence tell us the shape of what happens at death; they do not tell us what happened at any particular death. Purgatory is a doctrine about a state, not a report on the whereabouts of anyone. And the fact that the Church names the blessed while making no corresponding declaration in the other direction is not a rhetorical courtesy: it is what the doctrine of judgment actually licenses, since sanctity can be established from the visible evidence of a life and its effects, while damnation would require exactly the interior verdict that has been reserved.\endnote{The observation in this sentence about the asymmetry between the Church's naming of the blessed and her silence about the damned is this article's own and is labelled as project synthesis in the terminal appendix. This article did not examine the canonical procedures of canonization and makes no claim about their juridical form or their doctrinal note; the point made here is only that the doctrine of judgment set out in the preceding sections licenses the asymmetry, not that the asymmetry is itself a defined teaching. Readers should note that the Church's teaching on hell, and on what may and may not be said about who is there, is the subject of the companion volume in this series and is not developed here.}

\subsection{The end of the matter}

Set the whole file down and the doctrine that emerges is smaller than its reputation and harder than its caricature.

It is smaller because the defined content fits on a page: an immediate and irrevocable retribution at death, purification interposed where needed, the vision of God before the resurrection, a general resurrection of all in their own bodies, and an appearance of all before Christ's tribunal on a day nobody knows. Everything else in the received picture---the mechanics of the two hearings, the geography, the fire's physics, the duration, the courtroom procedure---is theology of varying quality, some of it excellent, none of it defined, and much of it fenced off from popular preaching by the one council that legislated on the subject and by the most recent Roman instruction on it.

It is harder because the two things that \emph{are} defined are the two that a modern reader would most like to soften. The first sentence is passed at death and it is final; there is no purgatorial second chance, because purgatory is for the saved. And the second judgment is public. The scriptural vocabulary for it is unrelenting on this point: made manifest, brought to light, laid bare, revealed, published in the light, preached on the housetops. Whatever consolation the doctrine offers---and Augustine thought the consolation immense, since it is the end of the ignorant questioning why the wicked prosper---it is not the consolation of privacy.

And that, finally, is why the second judgment is not redundant. A man's own case closes when he dies. What he did does not close when he dies; it goes on working in a world he has left, in memories that get him wrong, in children who are not like him, in consequences he never saw, in a body that is still being done things to. Only when history stops is there anything like the whole of him to look at. The tradition put the two judgments where it did because it believed both of two things that are hard to believe together: that a man is answerable now, without extension or appeal; and that a man is more than the instant of his answering. The last day is not a rehearing. It is the moment at which the sentence already passed on a person becomes the truth visible about a life---the whole of it, including the part that outlived him.
